\section{Identificación y descripción del problema}

El proceso de secado del café constituye una de las etapas críticas en la cadena de valor del sector cafetero colombiano, al influir directamente en la calidad física, química y sensorial del grano. Sin embargo, en la actualidad, este proceso se realiza mayoritariamente mediante métodos tradicionales basados en el control de temperatura y condiciones ambientales variables, lo cual limita la capacidad de controlar de manera precisa la cinética de secado y genera una alta variabilidad en la calidad del producto final.

Esta problemática se acentúa en el contexto de los cafés diferenciados, como los procesos honey y natural, los cuales requieren condiciones de secado más controladas debido a su mayor contenido de mucílago y su sensibilidad a fenómenos como la fermentación no controlada, la aparición de defectos y la pérdida de atributos sensoriales. A pesar del creciente interés en estos cafés por parte de mercados especializados, el sector cafetero colombiano carece de tecnologías que permitan estandarizar y escalar estos procesos a nivel industrial, lo que limita su competitividad frente a otros países productores.

Adicionalmente, a nivel de finca, muchos caficultores no cuentan con la infraestructura ni el conocimiento técnico necesarios para realizar un secado adecuado, lo que se traduce en pérdidas económicas asociadas a la disminución de la calidad del café. Esta situación evidencia una brecha significativa entre las condiciones reales de producción y los estándares requeridos por los mercados internacionales, especialmente en el segmento de cafés especiales.

En respuesta a estas limitaciones, se han promovido modelos de transformación del sistema productivo, como el esquema de compra de café en cereza (CERPER), impulsado por Cenicafé y la Federación Nacional de Cafeteros, el cual plantea la centralización del procesamiento en plantas de beneficio especializadas. No obstante, la implementación de este tipo de modelos requiere el desarrollo de tecnologías eficientes y escalables que garanticen la calidad del producto y la sostenibilidad del proceso.

En este contexto, se identifica una brecha tecnológica en el desarrollo de sistemas de secado que permitan un control preciso de las variables críticas del proceso, particularmente la humedad relativa, la cual juega un papel determinante en la transferencia de masa durante el secado del café. Los enfoques tradicionales centrados en la temperatura no permiten un control adecuado de esta variable, limitando la eficiencia del proceso y la calidad del producto final.

Por lo tanto, surge la necesidad de desarrollar y escalar tecnologías de secado innovadoras que permitan superar estas limitaciones, facilitando la estandarización de procesos de cafés diferenciados, la mejora en la calidad del producto, la reducción de pérdidas en la cadena productiva y la generación de mayores ingresos para los caficultores. Asimismo, dichas tecnologías deben ser compatibles con esquemas de producción sostenibles y adaptables a modelos de procesamiento centralizado, contribuyendo a la modernización del sector cafetero colombiano.